%%
%% part2-book-components/ch14-bibliography.tex
%%
%% Chapter 14: Bibliography — biblatex + biber, citations,
%% and the bibliography display.
%%

\chapter{منابع و کتاب‌نامه}
\label{chap:bibliography}

\section{چرا این موضوع مهم است؟}
\label{sec:bibliography-why}

استناددهی یکی از ارکان اساسی نگارش علمی است.
هر ادعای علمی باید به منبع معتبر ارجاع داده شود
و خواننده باید بتواند به‌راحتی منابع را پیدا
کند.

\lr{MatinBook} از \lr{biblatex} با موتور
\lr{biber} برای مدیریت منابع استفاده می‌کند.
این ترکیب، استاندارد مدرن \lr{LaTeX} است و از
تمام قابلیت‌های پیشرفته‌ی کتاب‌نامه پشتیبانی
می‌کند.

\mnote{موتور \lr{biber} جایگزین \lr{bibtex} قدیمی است و از \lr{Unicode} و فارسی پشتیبانی می‌کند.}

\section{فایل \lr{.bib}}
\label{sec:bibliography-bibfile}

منابع در یک فایل جداگانه با پسوند \file{.bib}
ذخیره می‌شوند. هر منبع، یک ورودی \lr{BibTeX}
دارد:

\begin{latin}
\begin{minted}{bibtex}
@book{knuth1984,
    author    = {Donald E. Knuth},
    title     = {The {\TeX}book},
    year      = {1984},
    publisher = {Addison-Wesley},
    address   = {Reading, MA}
}

@article{einstein1905,
    author  = {Albert Einstein},
    title   = {On the Electrodynamics of Moving Bodies},
    journal = {Annalen der Physik},
    volume  = {17},
    pages   = {891--921},
    year    = {1905}
}
\end{minted}
\end{latin}

\mnote{کلید هر منبع (مثل \lr{knuth1984}) باید یکتا باشد. برای استناد، از همین کلید استفاده می‌کنید.}

\section{بارگذاری فایل \lr{.bib}}
\label{sec:bibliography-loading}

برای بارگذاری فایل \lr{.bib}، از دستور
\cmd{addbibresource} در پیش از
\cmd{begin\{document\}} استفاده کنید:

\begin{latin}
\begin{minted}{latex}
\documentclass{matinbook}

\addbibresource{references.bib}

\begin{document}
...
\end{document}
\end{minted}
\end{latin}

\mnote{فایل \lr{.bib} باید در همان پوشه‌ی \lr{main.tex} یا در مسیر مشخص‌شده باشد.}

\section{استناد}
\label{sec:bibliography-citation}

\lr{biblatex} چهار دستور اصلی برای استناد ارائه
می‌دهد:

\begin{itemize}
    \item \cmd{cite\{key\}}: استناد ساده.
    \item \cmd{parencite\{key\}}: استناد پرانتزی.
    \item \cmd{textcite\{key\}}: استناد متنی.
    \item \cmd{autocite\{key\}}: استناد خودکار.
\end{itemize}

\subsection{استناد ساده}
\label{sec:bibliography-cite}

\begin{latin}
\begin{minted}{latex}
|\pc{کتاب}|\cite{knuth1984}|\pc{مرجع استاندارد}|\lr{\TeX}|\pc{است.}|%
\end{minted}
\end{latin}

خروجی: کتاب \cite{knuth1984} مرجع استاندارد
\lr{\TeX} است.

\mnote{دستور \lr{\textbackslash cite} شماره‌ی منبع را در براکت نمایش می‌دهد: \lr{[1]}.}

\subsection{استناد پرانتزی}
\label{sec:bibliography-parencite}

\begin{latin}
\begin{minted}{latex}
|\pc{این الگوریتم در}|\parencite{cormen2009}|\pc{معرفی شده است.}|%
\end{minted}
\end{latin}

خروجی: این الگوریتم در \parencite{cormen2009}
معرفی شده است.

\mnote{در سبک عددی، \lr{\textbackslash parencite} و \lr{\textbackslash cite} یکسان نمایش داده می‌شوند.}

\subsection{استناد متنی}
\label{sec:bibliography-textcite}

\begin{latin}
\begin{minted}{latex}
\textcite{cormen2009}|\pc{این الگوریتم را معرفی کرد.}|%
\end{minted}
\end{latin}

خروجی: \textcite{cormen2009} این الگوریتم را
معرفی کرد.

\mnote{دستور \lr{\textbackslash textcite} نام نویسنده را در متن و شماره‌ی منبع را در براکت نمایش می‌دهد.}

\subsection{استناد به چند منبع}
\label{sec:bibliography-multiple}

\begin{latin}
\begin{minted}{latex}
|\pc{منابع}|\cite{knuth1984,lamport1994}|\pc{برای مطالعه توصیه می‌شوند.}|%
\end{minted}
\end{latin}

خروجی: منابع \cite{knuth1984,lamport1994} برای
مطالعه توصیه می‌شوند.

\mnote{برای استناد به چند منبع، کلیدها را با کاما جدا کنید.}

\section{نمایش کتاب‌نامه}
\label{sec:bibliography-print}

برای نمایش کتاب‌نامه، از دستور
\cmd{printbibliography} در \lr{backmatter}
استفاده کنید:

\begin{latin}
\begin{minted}{latex}
\backmatter
\frontmattergeometry

\printbibliography[title={|\pc{منابع و مراجع}|}]
\end{minted}
\end{latin}

\mnote{کتاب‌نامه را \textbf{هرگز} داخل \lr{\textbackslash begin\{latin\}} قرار ندهید. \lr{xepersian} جهت متن را به‌درستی مدیریت می‌کند.}

\subsection{عنوان فارسی}
\label{sec:bibliography-persian-title}

عنوان کتاب‌نامه را با آرگومان \lr{title} تعیین
کنید:

\begin{latin}
\begin{minted}{latex}
\printbibliography[title={|\pc{منابع و مراجع}|}]
\end{minted}
\end{latin}

\mnote{عنوان پیش‌فرض کتاب‌نامه «منابع» است. با آرگومان \lr{title} می‌توانید آن را تغییر دهید.}

\section{سبک کتاب‌نامه}
\label{sec:bibliography-style}

\lr{MatinBook} از سبک \textbf{عددی} (\lr{numeric})
استفاده می‌کند. یعنی استنادها به‌صورت شماره
نمایش داده می‌شوند: \lr{[1]}، \lr{[2]} و...

\begin{latin}
\begin{minted}{latex}
\RequirePackage[
    backend=biber,
    style=numeric,
    sorting=none
]{biblatex}
\end{minted}
\end{latin}

\mnote{ترتیب \lr{sorting=none} یعنی منابع به ترتیب استناد در کتاب‌نامه ظاهر می‌شوند، نه به ترتیب الفبایی.}

\subsection{سبک‌های جایگزین}
\label{sec:bibliography-alternative-styles}

اگر سبک دیگری می‌خواهید، می‌توانید آن را در
فایل \file{mb-core.sty} تغییر دهید:

\begin{itemize}
    \item \lr{style=numeric}: عددی (پیش‌فرض).
    \item \lr{style=alphabetic}: الفبایی-عددی
    (مثل \lr{[Knu84]}).
    \item \lr{style=authoryear}: نویسنده-سال
    (مثل \lr{(Knuth, 1984)}).
\end{itemize}

\mnote{برای کتاب‌های فنی، سبک عددی توصیه می‌شود چون فشرده‌تر است.}

\section{انواع ورودی \lr{BibTeX}}
\label{sec:bibliography-entry-types}

جدول \cref{tab:bib-types} انواع اصلی ورودی
\lr{BibTeX} را نشان می‌دهد.

\begin{matintable}{انواع ورودی \lr{BibTeX}}{tab:bib-types}
\begin{tblr}{
    width = \textwidth,
    colspec = {l X[l]},
    hlines, vlines,
    colsep = 8pt,
    rowsep = 4pt,
    row{1} = {font=\bfseries},
}
نوع & کاربرد \\
\lr{@book} & کتاب \\
\lr{@article} & مقاله‌ی ژورنالی \\
\lr{@inproceedings} & مقاله‌ی کنفرانسی \\
\lr{@online} & وب‌سایت \\
\lr{@phdthesis} & پایان‌نامه‌ی دکتری \\
\lr{@mastersthesis} & پایان‌نامه‌ی کارشناسی ارشد \\
\lr{@techreport} & گزارش فنی \\
\lr{@misc} & سایر \\
\end{tblr}
\end{matintable}

\mnote{هر نوع ورودی، فیلدهای مخصوص خود را دارد. برای اطلاعات کامل، به مستندات \lr{biblatex} مراجعه کنید.}

\section{منابع فارسی}
\label{sec:bibliography-persian}

\lr{biblatex} با \lr{biber} از منابع فارسی
پشتیبانی می‌کند. برای یک منبع فارسی، از فیلد
\lr{langid} استفاده کنید:

\begin{latin}
\begin{minted}{bibtex}
@book{persian-book,
    author    = {|\pc{نام نویسنده}|},
    title     = {|\pc{عنوان کتاب}|},
    year      = {|\pnum{۱۴۰۴}|},
    publisher = {|\pc{نام ناشر}|},
    langid    = {persian}
}
\end{minted}
\end{latin}

\mnote{فیلد \lr{langid = \{persian\}} باعث می‌شود \lr{biblatex} جهت متن را راست‌به‌چپ کند.}

\section{چرخه‌ی کامپایل}
\label{sec:bibliography-compile}

برای ساخت کتاب‌نامه، باید چرخه‌ی چهار مرحله‌ای
را اجرا کنید:

\begin{latin}
\begin{minted}{bash}
# |\pc{مرحله‌ی اول: کامپایل اول}|%
xelatex -shell-escape main.tex

# |\pc{مرحله‌ی دوم: پردازش کتاب‌نامه}|%
biber main

# |\pc{مرحله‌ی سوم: کامپایل دوم}|%
xelatex -shell-escape main.tex

# |\pc{مرحله‌ی چهارم: کامپایل سوم}|%
xelatex -shell-escape main.tex
\end{minted}
\end{latin}

\mnote{اگر یکی از این مراحل را فراموش کنید، کتاب‌نامه خالی خواهد بود یا استنادها به‌صورت «؟؟» نمایش داده می‌شوند.}

\section{استناد به منابع بدون استناد}
\label{sec:bibliography-nocite}

اگر می‌خواهید منبعی در کتاب‌نامه ظاهر شود
بدون اینکه در متن به آن استناد کنید، از دستور
\cmd{nocite} استفاده کنید:

\begin{latin}
\begin{minted}{latex}
\nocite{knuth1984}
\nocite{*}
\end{minted}
\end{latin}

\begin{itemize}
    \item \cmd{nocite\{key\}}: یک منبع خاص.
    \item \cmd{nocite\{*\}}: همه‌ی منابع فایل
    \lr{.bib}.
\end{itemize}

\mnote{استفاده از \lr{\textbackslash nocite\{*\}} برای کتاب‌های مرجع مفید است، چون همه‌ی منابع را در کتاب‌نامه می‌آورد.}

\section{خطاهای رایج}
\label{sec:bibliography-mistakes}

\subsection{فراموش کردن \cmd{addbibresource}}
\label{sec:bibliography-mistake-addbibresource}

\textbf{مشکل:} اگر \cmd{addbibresource} را فراموش
کنید، کتاب‌نامه خالی خواهد بود.

\textbf{راه‌حل:} همیشه \cmd{addbibresource} را
در پیش از \cmd{begin\{document\}} قرار دهید.

\subsection{فراموش کردن \lr{biber}}
\label{sec:bibliography-mistake-biber}

\textbf{مشکل:} اگر \lr{biber} را اجرا نکنید،
استنادها به‌صورت «؟؟» نمایش داده می‌شوند.

\textbf{راه‌حل:} چرخه‌ی کامل چهار مرحله‌ای را
اجرا کنید.

\subsection{قرار دادن کتاب‌نامه در \env{latin}}
\label{sec:bibliography-mistake-latin}

\textbf{مشکل:} اگر \cmd{printbibliography} را
داخل \cmd{begin\{latin\}} قرار دهید، عنوان
فارسی «منابع و مراجع» به‌هم می‌ریزد.

\textbf{راه‌حل:} کتاب‌نامه را \textbf{هرگز} داخل
\env{latin} قرار ندهید.

\subsection{کلید تکراری}
\label{sec:bibliography-mistake-duplicate}

\textbf{مشکل:} اگر دو منبع کلید یکسان داشته
باشند، \lr{biber} خطا می‌دهد.

\textbf{راه‌حل:} کلیدهای یکتا برای هر منبع
تعریف کنید.

\section{تمرین‌ها}
\label{sec:bibliography-exercises}

\begin{exercise}
    سه منبع (یک کتاب، یک مقاله، یک وب‌سایت)
    در فایل \lr{.bib} تعریف کنید.
    \label{exc:bibliography-bibfile}
\end{exercise}

\begin{exercise}
    به هر سه منبع استناد کنید و کتاب‌نامه را
    نمایش دهید.
    \label{exc:bibliography-cite}
\end{exercise}

\begin{exercise}
    یک منبع فارسی با فیلد \lr{langid} تعریف
    کنید و به آن استناد دهید.
    \label{exc:bibliography-persian}
\end{exercise}

\begin{exercise}
    از \cmd{nocite\{*\}} استفاده کنید و همه‌ی
    منابع را در کتاب‌نامه نمایش دهید.
    \label{exc:bibliography-nocite}
\end{exercise}

\section{خلاصه}
\label{sec:bibliography-summary}

در این فصل، با منابع و کتاب‌نامه در \lr{MatinBook}
آشنا شدیم:

\begin{itemize}
    \item \lr{MatinBook} از \lr{biblatex} با
    موتور \lr{biber} استفاده می‌کند.

    \item منابع در فایل \lr{.bib} ذخیره می‌شوند.

    \item فایل \lr{.bib} با \cmd{addbibresource}
    بارگذاری می‌شود.

    \item چهار دستور استناد: \cmd{cite}،
    \cmd{parencite}، \cmd{textcite}، \cmd{autocite}.

    \item کتاب‌نامه با \cmd{printbibliography}
    نمایش داده می‌شود.

    \item سبک پیش‌فرض، عددی (\lr{numeric}) است.

    \item منابع فارسی با \lr{langid = \{persian\}}
    تعریف می‌شوند.

    \item چرخه‌ی کامپایل چهار مرحله‌ای: \lr{xelatex}
    + \lr{biber} + \lr{xelatex} + \lr{xelatex}.

    \item خطاهای رایج شامل فراموش کردن
    \cmd{addbibresource}، فراموش کردن \lr{biber}،
    قرار دادن کتاب‌نامه در \env{latin}، و کلید
    تکراری است.
\end{itemize}

در فصل بعدی (\cref{chap:index})، با نمایه در
\lr{MatinBook} آشنا می‌شوید و یاد می‌گیرید
که چگونه یک نمایه‌ی حرفه‌ای بسازید.